% Exercices type bac — Dérivation v17
% À intégrer dans le document Dérivation Première spécialité.

\section*{Sujets type bac — Dérivation}

\subsection*{Sujet 1 — Courbe, signe et variations d'une fonction rationnelle}
On considère la fonction $g$ définie sur $\mathbb{R}$ par
\[
  g(x)=x^2-5x+4.
\]
On considère aussi la fonction $f$ définie sur l'intervalle $[1;8]$ par
\[
  f(x)=x-5+\frac4x.
\]
\begin{enumerate}
  \item Factoriser $g(x)$, puis étudier le signe de $g(x)$ sur $\mathbb{R}$.
  \item Vérifier que, pour tout $x\in[1;8]$, $f(x)=\dfrac{g(x)}{x}$.
  \item En déduire le signe de $f(x)$ sur $[1;8]$.
  \item Calculer $f'(x)$.
  \item Montrer que $f'(x)=\dfrac{(x-2)(x+2)}{x^2}$.
  \item Étudier le signe de $f'(x)$ sur $[1;8]$.
  \item Dresser le tableau de variations de $f$ sur $[1;8]$.
  \item Déterminer le minimum de $f$ sur $[1;8]$.
  \item Interpréter graphiquement le résultat obtenu à la question 3.
\end{enumerate}

\subsection*{Sujet 2 — Fonction avec paramètre et tangente}
On considère la fonction $f$ définie sur $\mathbb{R}$ par
\[
  f(x)=(ax+b)e^{-x},
\]
où $a$ et $b$ sont deux réels. Sa courbe représentative passe par le point $A(0;4)$. La tangente à la courbe au point $A$ passe par le point $B(1;7)$.
\begin{enumerate}
  \item En utilisant $f(0)=4$, déterminer $b$.
  \item Calculer le coefficient directeur de la tangente en $A$.
  \item En déduire $f'(0)$.
  \item Montrer que, pour tout réel $x$,
  \[
    f'(x)=(a-ax-b)e^{-x}.
  \]
  \item En déduire la valeur de $a$, puis l'expression de $f(x)$.
  \item Étudier le signe de $f'(x)$.
  \item Dresser le tableau de variations de $f$ sur $\mathbb{R}$.
  \item Déterminer le maximum de $f$ sur $\mathbb{R}$.
  \item Donner une équation de la tangente à la courbe au point $A$.
\end{enumerate}

\subsection*{Sujet 3 — Concentration et seuil réglementaire}
Après injection d'un produit, la concentration dans le sang est modélisée, en mg/L, par la fonction $C$ définie sur $[0;5]$ par
\[
  C(t)=4t e^{-t}.
\]
\begin{enumerate}
  \item Calculer $C(0)$.
  \item Montrer que, pour tout $t\in[0;5]$,
  \[
    C'(t)=4(1-t)e^{-t}.
  \]
  \item Étudier le signe de $C'(t)$ sur $[0;5]$.
  \item Dresser le tableau de variations de $C$ sur $[0;5]$.
  \item Déterminer l'instant où la concentration est maximale.
  \item Donner la valeur exacte de cette concentration maximale.
  \item La réglementation impose une concentration inférieure ou égale à $1$ mg/L. Justifier, à l'aide du tableau de variations, qu'il existe un intervalle de temps durant lequel le sportif ne respecte pas la réglementation.
  \item On admet que $e^2>4$. Montrer que $C(2)<1$.
  \item En déduire que la concentration repasse sous le seuil de $1$ mg/L entre $1$ h et $2$ h après l'injection.
\end{enumerate}
